\documentclass[10pt]{article}
\usepackage[landscape, margin=0.3in]{geometry}
\usepackage{multicol}
\usepackage{amsmath}
\usepackage{enumitem}
\setlist{nosep}
\pagenumbering{gobble}

\begin{document}
\begin{multicols}{4}

%========================
\section*{Geology Exam 2}
%========================

\textbf{Half-Life (Radiometric Dating):}
Time required for half of a radioactive parent isotope to decay into daughter product.
\begin{itemize}
\item 50\% parent = 1 half-life (age = 1 $\times$ HL)
\item 25\% parent = 2 half-lives
\item 12.5\% parent = 3 half-lives
\item 75\% parent = 0.5 half-life (e.g., HL = 200 yrs $\rightarrow$ age = 100 yrs)
\item Test logic: More daughter product = older sample
\end{itemize}

\textbf{Earthquake Magnitude:}
Logarithmic scale measuring energy released.
\begin{itemize}
\item +1 magnitude = 32$\times$ more energy
\item +2 magnitude = 1024$\times$ more energy
\item Test logic: small magnitude increase = huge energy difference
\end{itemize}

\textbf{Stream Discharge:}
Volume of water flowing per unit time.
\[
Q = A \times v
\]
\begin{itemize}
\item A = cross-sectional area, v = velocity
\item Discharge increases downstream as tributaries add water
\item Test logic: widest/deepest part of river $\rightarrow$ highest discharge
\end{itemize}

%========================
\section*{Key Orders / Rankings}
%========================

\textbf{Geologic Time Scale (Largest $\rightarrow$ Smallest):}
\begin{itemize}
\item Eon: largest division (billions of years)
\item Era: subdivisions of eons
\item Period: subdivisions of eras
\item Epoch: smallest division (most recent detail)
\end{itemize}

\textbf{Seismic Wave Arrival Order:}
\begin{itemize}
\item P-waves: fastest, compressional, arrive first
\item S-waves: slower, shear motion, arrive second
\item Surface waves: slowest, largest amplitude, cause most damage
\item Test logic: biggest shaking = last waves
\end{itemize}

\textbf{Chemical Weathering Rate:}
\begin{itemize}
\item Hot/Wet: fastest (water + heat accelerate reactions)
\item Temperate: moderate
\item Cold/Dry: slowest (little water, low reaction rates)
\end{itemize}

%========================
\section*{Weathering}
%========================

\textbf{Weathering:}
Breakdown of rock at Earth’s surface without movement.

\textbf{Mechanical Weathering:}
Physical breakdown into smaller pieces without chemical change.
\begin{itemize}
\item Frost Wedging: water freezes in cracks, expands, and splits rock
\item Salt Crystal Growth: salt expands in pores, breaking rock (common in deserts/coasts)
\item Unloading/Sheeting: removal of overlying rock reduces pressure, causing outer layers to peel (exfoliation)
\item Biological Activity: roots grow into cracks and expand them
\item Hydration: repeated wetting/drying causes expansion and cracking (mud cracks)
\end{itemize}

\textbf{Weathering Geometries (What rocks look like):}
\begin{itemize}
\item Granular Disintegration: rock crumbles into individual grains (granite)
\item Spheroidal Weathering: corners weather faster, producing rounded shapes
\item Shattering: rock breaks into jagged, irregular fragments
\end{itemize}

\textbf{Chemical Weathering:}
Breakdown by chemical reactions that change mineral composition.
\begin{itemize}
\item Dissolution: minerals dissolve completely in water (halite)
\item Carbonation: carbonic acid reacts with calcite → caves/karst
\item Oxidation: iron reacts with oxygen → rust (reddish color)
\item Hydrolysis: silicate minerals (feldspar) convert into clay
\end{itemize}

\textbf{Controls on Weathering:}
\begin{itemize}
\item Climate: hot/wet = chemical; cold/dry = mechanical
\item Rock type: some minerals weather faster (feldspar > quartz)
\end{itemize}

\textbf{Regolith:}
Layer of loose rock and mineral fragments produced by weathering covering Earth’s surface.

%========================
\section*{Soils}
%========================

\textbf{Soil:}
Mixture of weathered rock, organic material, water, and air that supports plant life.

\textbf{Soil Formation Factors:}
\begin{itemize}
\item Parent Material: original rock type
\item Climate: most important (controls weathering rate)
\item Organisms: add organic matter and mix soil
\item Time: longer time = thicker, more developed soil
\item Topography: slope affects erosion and drainage
\end{itemize}

\textbf{Soil Horizons (Profile):}
\begin{itemize}
\item O: organic material (leaves, decomposed matter)
\item A: topsoil; nutrient-rich, supports plant growth
\item E: leaching zone; minerals washed out
\item B: accumulation zone; minerals deposited from above
\item C: partially weathered parent material
\end{itemize}

\textbf{Soil Texture (Particle Size):}
\begin{itemize}
\item Clay ($<0.002$ mm): very small, holds water, low permeability
\item Silt (0.002–0.06 mm): smooth, moderate properties
\item Sand (0.06–2 mm): large, drains quickly
\end{itemize}

\textbf{Soil Structure (Shape of aggregates):}
\begin{itemize}
\item Platy: flat layers, restrict water movement
\item Prismatic: vertical columns, moderate drainage
\item Blocky: cube-like, good for roots
\item Spheroidal: rounded, high permeability
\end{itemize}

\textbf{Bioturbation:}
Mixing of soil by organisms (worms, roots), improving aeration and nutrient distribution.

\textbf{Soil Origin:}
\begin{itemize}
\item Residual: forms directly from underlying rock
\item Transported: formed from sediment moved by water, wind, or ice
\end{itemize}

%========================
\section*{Sedimentary Rocks}
%========================

\textbf{Sedimentary Rocks:}
Rocks formed from the accumulation, compaction, and cementation of sediments or from chemical precipitation.

\textbf{Clastic Sediments (Broken Rock Fragments):}
Classified by grain size.
\begin{itemize}
\item Gravel ($>2$ mm): large fragments
\item Sand ($1/16 - 2$ mm): visible grains
\item Mud ($<1/16$ mm): very fine particles (clay + silt)
\end{itemize}

\textbf{Transport Trends (Distance from Source):}
\begin{itemize}
\item Grain size decreases (large $\rightarrow$ small)
\item Roundness increases (angular $\rightarrow$ smooth)
\item Sorting improves (mixed sizes $\rightarrow$ uniform)
\item Test logic: well-rounded + well-sorted = long transport
\end{itemize}

\textbf{Common Clastic Rocks:}
\begin{itemize}
\item Breccia: angular gravel; indicates little transport
\item Conglomerate: rounded gravel; indicates water transport
\item Sandstone: sand-sized grains; common in beaches/deserts
\item Shale: mud-sized particles; smooth and fissile
\end{itemize}

\textbf{Sandstone Types:}
\begin{itemize}
\item Arkose: feldspar-rich, angular grains → short transport from granite source
\item Graywacke: poorly sorted with matrix → rapid deposition (e.g., underwater landslides)
\end{itemize}

\textbf{Shale Properties:}
\begin{itemize}
\item Fissility: ability to split into thin sheets
\item Laminae: very thin sediment layers visible in rock
\item Test logic: only sedimentary rock that splits easily into layers
\end{itemize}

\textbf{Chemical / Biochemical Rocks:}
Form from dissolved minerals or biological processes.
\begin{itemize}
\item Limestone: calcite (CaCO$_3$); fizzes with HCl acid (key ID test)
\item Coquina: loosely cemented shell fragments
\item Chalk: microscopic marine organisms
\item Travertine: forms in caves/springs from mineral precipitation
\item Dolomite: Mg-rich carbonate; reacts weakly unless powdered
\end{itemize}

\textbf{Evaporites:}
Rocks formed when water evaporates, leaving minerals behind.
\begin{itemize}
\item Halite (rock salt), Gypsum
\end{itemize}

\textbf{Coal Formation (Organic Sedimentary):}
\begin{itemize}
\item Peat $\rightarrow$ Lignite $\rightarrow$ Bituminous $\rightarrow$ Anthracite
\item Forms in swamps where stagnant water slows bacterial decay, preserving plant material
\end{itemize}

\textbf{Carbon Cycle (Sedimentary Role):}
\begin{itemize}
\item Limestone = largest carbon reservoir on Earth
\item Sequestration: photosynthesis, burial of organic matter, limestone formation
\item Release: burning fossil fuels, volcanic eruptions
\end{itemize}

\textbf{Sedimentary Structures (Clues to Environment):}
\begin{itemize}
\item Cross-bedding: angled layers formed by wind/water currents
\begin{itemize}
\item Flow direction = toward steepest side (down-current)
\end{itemize}
\item Graded bedding: coarse at bottom, fine at top → rapid deposition (turbidity currents)
\item Ripple marks:
\begin{itemize}
\item Oscillation: symmetrical → wave motion (back-and-forth)
\item Current: asymmetrical → one-direction flow
\end{itemize}
\end{itemize}

\textbf{Sedimentary Facies:}
Set of rock characteristics that represent a specific environment.
\begin{itemize}
\item Facies change laterally (e.g., beach sand → offshore mud)
\end{itemize}

\textbf{Depositional Environments:}
\begin{itemize}
\item Alluvial Fan: coarse sediment (breccia/conglomerate)
\item Delta: sediment deposited where river enters water body
\item Swamp: organic-rich → coal formation
\item Glacial: poorly sorted till (mixed sizes)
\end{itemize}

%========================
\section*{Metamorphism}
%========================

\textbf{Metamorphic Rocks:}
Rocks altered by heat, pressure, and chemically active fluids without melting.

\textbf{Foliation:}
Alignment of minerals into layers due to directed pressure.
\begin{itemize}
\item Slate: lowest grade, very fine layers
\item Phyllite: silky sheen, slightly larger crystals
\item Schist: visible mineral grains (mica)
\item Gneiss: high-grade, banded light/dark minerals
\item Test logic: increasing grade = increasing crystal size and separation
\end{itemize}

\textbf{Non-Foliated Rocks:}
Do not show layering; form under uniform pressure.
\begin{itemize}
\item Marble: from limestone; reacts with acid
\item Quartzite: from sandstone; very hard, scratches glass
\end{itemize}

\textbf{Protolith:}
Original rock before metamorphism.
\begin{itemize}
\item Limestone $\rightarrow$ Marble
\item Sandstone $\rightarrow$ Quartzite
\end{itemize}

\textbf{Orogeny:}
Process of mountain building caused by tectonic plate collision.

\textbf{Example:}
\begin{itemize}
\item Appalachian Mountains formed from collision of Africa and North America
\end{itemize}

%========================
\section*{Geologic Time}
%========================

\textbf{Absolute Dating:}
Determines numeric age using radioactive decay.

\textbf{Relative Dating Principles:}
\begin{itemize}
\item Superposition: oldest layers on bottom, youngest on top
\item Cross-cutting: feature that cuts another is younger
\item Inclusions (Xenoliths): included rock is older than host
\end{itemize}

\textbf{Unconformities (Gaps in Time):}
\begin{itemize}
\item Angular: tilted layers below flat layers → deformation before deposition
\item Disconformity: parallel layers with missing time → erosion/non-deposition
\item Nonconformity: igneous/metamorphic below sedimentary
\item Test logic: unconformity = missing geologic record
\end{itemize}

\textbf{Index Fossils:}
Used to correlate rock layers.
\begin{itemize}
\item Must be widespread geographically
\item Must exist for short geologic time
\end{itemize}

\textbf{Fossil Types:}
\begin{itemize}
\item Permineralization: minerals fill pores, preserving structure
\item Trace fossils: footprints, burrows, coprolites
\end{itemize}

\textbf{Key Dates:}
\begin{itemize}
\item Earth: 4.56 billion years
\item Oldest mineral (zircon): 4.4 bya
\item Oldest rock: 4.013 bya
\item Cambrian Explosion: 540 million years ago
\end{itemize}

\textbf{Hadean Eon:}
Earliest Earth history (>4 bya); extremely hot, rocks frequently remelted

\textbf{Isotopes:}
Atoms of same element (same protons) with different numbers of neutrons.
\begin{itemize}
\item Used in radiometric dating
\end{itemize}

%========================
\section*{Earthquakes}
%========================

\textbf{Earthquake:}
Sudden release of stored elastic energy in rocks along a fault, causing ground shaking.

\textbf{Elastic Rebound Theory:}
Rocks deform elastically as stress builds, then suddenly break and release energy.
\begin{itemize}
\item Test logic: energy is stored BEFORE the earthquake, released DURING rupture
\end{itemize}

\textbf{Seismic Waves:}
Energy waves traveling through Earth.
\begin{itemize}
\item P-waves (Primary): compressional motion, travel through solids/liquids, fastest → arrive first
\item S-waves (Secondary): shear motion, only through solids → arrive second
\item Surface waves: travel along surface, slowest but largest amplitude → cause most damage
\end{itemize}

\textbf{Epicenter Location:}
Point on Earth's surface directly above focus.
\begin{itemize}
\item Found using at least 3 seismic stations measuring P–S arrival time differences
\end{itemize}

\textbf{Magnitude vs Intensity:}
\begin{itemize}
\item Magnitude: total energy released (objective, one value per earthquake)
\item Intensity: observed shaking/damage (varies by location)
\end{itemize}

\textbf{Earthquake Sequence:}
\begin{itemize}
\item Foreshocks: smaller quakes before main event
\item Mainshock: largest energy release
\item Aftershocks: smaller quakes after, as crust readjusts
\end{itemize}

\textbf{Hazards:}
\begin{itemize}
\item Liquefaction: water-saturated sediment behaves like liquid → buildings sink
\item Tsunami: large sea wave from seafloor displacement
\end{itemize}

\textbf{New Madrid Seismic Zone:}
\begin{itemize}
\item Intraplate earthquakes (not at plate boundary)
\item Caused by reactivation of ancient faults from past supercontinent breakup
\end{itemize}

\textbf{Engineering for Safety:}
\begin{itemize}
\item Best material: wood (flexible, absorbs energy)
\item Worst: rigid materials (brick, concrete)
\item Best foundation: solid bedrock (least shaking)
\end{itemize}

%========================
\section*{Groundwater}
%========================

\textbf{Groundwater:}
Water stored beneath Earth's surface in pore spaces of rock/sediment.

\textbf{Porosity vs Permeability:}
\begin{itemize}
\item Porosity: amount of open space in material (storage capacity)
\item Permeability: ability of fluid to flow through material (connected pores)
\item Test logic: high porosity does NOT always mean high permeability
\end{itemize}

\textbf{Aquifer vs Aquitard:}
\begin{itemize}
\item Aquifer: permeable rock that stores and transmits water (sandstone)
\item Aquitard: impermeable layer that restricts flow (shale)
\end{itemize}

\textbf{Groundwater Flow:}
\begin{itemize}
\item Flows from high pressure to low pressure
\item In unsaturated zone: flows mostly straight downward
\item In saturated zone: follows curved paths based on pressure gradients
\item Described by Darcy’s Law (flow rate depends on permeability + pressure difference)
\end{itemize}

\textbf{Artesian System:}
\begin{itemize}
\item Confined aquifer between two aquitards
\item Water under pressure rises when tapped by a well
\item Requires recharge area at higher elevation
\end{itemize}

\textbf{Karst Topography:}
Landscape formed by dissolution of limestone.
\begin{itemize}
\item Features: sinkholes, caves, underground streams, natural bridges
\item Test logic: always associated with carbonate rock (limestone)
\end{itemize}

%========================
\section*{Surface Water}
%========================

\textbf{Stream Discharge (Review):}
Volume of water flowing per unit time.
\begin{itemize}
\item Increases downstream as tributaries add water
\end{itemize}

\textbf{Stream Features:}
\begin{itemize}
\item Meander: curved path of river
\item Cut Bank: outer curve; fastest flow → erosion
\item Point Bar: inner curve; slower flow → deposition
\item Test logic: erosion always on outside, deposition on inside
\end{itemize}

\textbf{Urbanization Effects:}
\begin{itemize}
\item Decreases lag time (water reaches stream faster)
\item Increases peak discharge (higher flood risk)
\item Caused by impervious surfaces (roads, buildings)
\end{itemize}

%========================
\section*{Mass Wasting}
%========================

\textbf{Mass Wasting:}
Downslope movement of material due to gravity.

\textbf{Driving Force:}
\begin{itemize}
\item Gravity (primary cause)
\end{itemize}

\textbf{Creep:}
\begin{itemize}
\item Slowest form of mass wasting
\item Movement occurs over long time scales
\item Identification clue: curved tree trunks, tilted fences
\item Test logic: if nothing “looks dramatic,” it is usually creep
\end{itemize}

%========================
\section*{Climate}
%========================

\textbf{Milankovitch Cycles:}
Variations in Earth's orbit affecting climate over long timescales.
\begin{itemize}
\item Eccentricity: shape of orbit (circular vs elliptical)
\item Obliquity: tilt of Earth's axis (affects seasons)
\item Precession: wobble of Earth's axis (changes timing of seasons)
\end{itemize}

\textbf{Albedo:}
Measure of how much solar radiation a surface reflects.
\begin{itemize}
\item High albedo (ice) → reflects more energy → cooling
\item Low albedo (water/land) → absorbs more → warming
\end{itemize}

\textbf{Albedo Feedback (Positive Feedback Loop):}
\begin{itemize}
\item Ice melts → albedo decreases → more heat absorbed → more melting
\end{itemize}

\textbf{Volcanic Effects:}
\begin{itemize}
\item SO$_2$ released into atmosphere forms aerosols
\item Increases albedo → reflects sunlight → cooling effect
\end{itemize}

\textbf{Sunspots (Course-Specific Logic):}
\begin{itemize}
\item Lower sunspot activity → more solar energy reaches Earth (per course)
\end{itemize}

\textbf{Climate Proxies:}
Indirect evidence used to reconstruct past climates.
\begin{itemize}
\item Ice cores: trap ancient air bubbles (atmospheric composition)
\item Oxygen isotopes:
\begin{itemize}
\item Lower $^{18}O/^{16}O$ ratio → warmer climate
\item Higher ratio → colder climate (more ice storage)
\end{itemize}
\end{itemize}

\end{multicols}
\end{document}